\documentclass[11pt]{article}
\usepackage[T1]{fontenc}
\usepackage[margin=1in]{geometry}
\usepackage{amsmath,amssymb,amsthm,mathtools,booktabs,array}
\usepackage[colorlinks=true,linkcolor=blue,urlcolor=blue,citecolor=blue]{hyperref}
\usepackage{microtype}
\hypersetup{pdftitle={A four dimensional finite volume criterion for a transfer gap},pdfauthor={Research note and proof audit},pdfsubject={A conditional lattice spectral criterion; not a continuum Yang--Mills construction}}
\newtheorem{theorem}{Theorem}[section]
\newtheorem{lemma}[theorem]{Lemma}
\newtheorem{proposition}[theorem]{Proposition}
\newtheorem{corollary}[theorem]{Corollary}
\theoremstyle{definition}
\newtheorem{assumption}[theorem]{Assumption}
\theoremstyle{remark}
\newtheorem{remark}[theorem]{Remark}
\newcommand{\Tr}{\operatorname{Tr}}
\newcommand{\spec}{\operatorname{spec}}
\newcommand{\Cov}{\operatorname{Cov}}
\newcommand{\norm}[1]{\left\lVert#1\right\rVert}
\newcommand{\ip}[2]{\langle#1,#2\rangle}
\newcommand{\one}{\mathbf 1}
\newcommand{\R}{\mathbb R}
\setlength{\emergencystretch}{3em}
\title{A four dimensional finite volume criterion for a transfer gap}
\author{Research note and proof audit}
\date{October 9, 2026}
\begin{document}
\maketitle
\begin{abstract}
We extend the elementary two-direction partition-function argument in the
repository's $O(4)$ manuscript to four directions. For a consistently
normalized rectangular-torus model admitting positive transfer traces in
every coordinate direction, define
$D(n)=16\log Z(n,n,n,n)-\log Z(2n,2n,2n,2n)$.
We prove that $D(n)\le 1/2048$ implies a transfer eigenvalue ratio at most
$2^{-1/n}$ on every spatial cube of side $2^k n$.
Under explicit local-observable and reflection-positivity hypotheses,
the same bound holds in every local limit along these dyadic tori.
The argument does not prove that the small-defect test holds on a weak-bare-coupling
continuum trajectory. It does not construct a continuum Yang--Mills theory,
prove nontriviality, or establish uniqueness of all thermodynamic states.
No novelty claim is made for the criterion.
\end{abstract}
\tableofcontents

\section{Scope and structural hypotheses}
This note concerns the finite-volume criterion of Family 215 in the supplied
repository, not its model-specific mass-generation and renormalization estimates.
The two-dimensional source is \cite[section ``A finite-volume criterion
for the full gap'']{repo-o4}; its use in the continuum argument appears in
\cite[section ``Trace estimates uniform in the cutoff'']{repo-o3}.
The repository revision inspected is
\texttt{fd4aeeb2ee4fc729c18d98444fed42fd0529eeeb}.

For an even side-length vector $x=(x_1,x_2,x_3,x_4)$ with all $x_i\ge4$,
let $Z(x)>0$ be a torus partition function and $F(x)=\log Z(x)$.
Let $D_i x$ double only coordinate $i$ and let $Dx=D_1D_2D_3D_4x$.
These operations commute. Write $\underline r=(r,r,r,r)$.

\begin{assumption}[Positive rectangular transfer traces]\label{ass:trace}
For every direction $i$ and fixed cross-section $\widehat x_i$ there is a
nonzero positive, self-adjoint, compact, trace-class operator
$K_{i,\widehat x_i}$ such that
\[
 Z(x)=\Tr K_{i,\widehat x_i}^{\,x_i}
\]
for every admitted $x_i$. Its largest eigenvalue is positive and simple.
The same function $Z$ and the same normalization are used in every direction.
\end{assumption}
Positive means positive semidefinite as an operator, not merely that an
integral kernel has positive entries. Only this assumption is needed for
the finite-cross-section theorem. Section~\ref{sec:local} states the
additional hypotheses for local limits.

Define directional and simultaneous doubling defects by
\begin{align}
 p_i(x)&=2F(x)-F(D_i x),\\
 D(r)&=16F(\underline r)-F(2\underline r).\label{eq:D}
\end{align}
Multiplying $Z(x)$ by $\exp(c\prod_i x_i)$ leaves every defect unchanged.
An arbitrary size-dependent or topology-dependent renormalization need not
do so, and need not preserve the transfer-trace hypothesis.

\begin{theorem}[Four dimensional finite-volume criterion]\label{thm:main}
Suppose Assumption~\ref{ass:trace} holds. Let $n\ge4$ be even. If
\[
 D(n)\le\varepsilon_*:=1/2048,
\]
then for every $k\ge0$, every coordinate direction, and the spatial cubic
cross-section of side $w_k=2^k n$, the first two transfer eigenvalues obey
\[
 \lambda_2(w_k)/\lambda_1(w_k)\le2^{-1/n}.
\]
If the operator has rank one, read $\lambda_2=0$. The full normalized
transfer contraction on the orthogonal complement of its leading
eigenvector has the same upper bound.
\end{theorem}
The assertion concerns dyadic cubic cross-sections, not all rectangles or
uniqueness of a Gibbs state. A lower bound on
$-\log(\lambda_2/\lambda_1)$ does not assert that it is finite.

\section{The one-direction spectral estimate}
\begin{lemma}[Trace concentration and squaring]\label{lem:squaring}
Let $K$ satisfy the operator hypotheses above, with eigenvalues
$\lambda_1>\lambda_2\ge\lambda_3\ge\cdots\ge0$. Put
\[
 p_m=2\log\Tr K^m-\log\Tr K^{2m},\qquad
 s_m=\sum_{j\ge2}(\lambda_j/\lambda_1)^m.
\]
Then $p_m\ge0$ and $s_m\le e^{p_m}-1$. If $p_m\le1/4$, then
\[
 (\lambda_2/\lambda_1)^m\le2p_m,\qquad p_{2m}\le4p_m^2.
\]
In particular $p_i(x)\ge0$ and
\begin{equation}
 p_i(D_i x)\le4p_i(x)^2\quad\text{if }p_i(x)\le1/4.\label{eq:square}
\end{equation}
\end{lemma}
\begin{proof}
Set $t_j=\lambda_j^m/\Tr K^m$. Then $\sum_jt_j=1$,
$t_1=(1+s_m)^{-1}$, and
\[
 e^{-p_m}=\sum_jt_j^2\le t_1\sum_jt_j=t_1\le1.
\]
For $0\le p_m\le1/4$, $e^{p_m}-1\le p_m e^{p_m}\le2p_m$.
Also $s_{2m}\le s_m^2$, so cancellation of the leading eigenvalue gives
\[
 p_{2m}=2\log(1+s_{2m})-\log(1+s_{4m})
 \le2s_{2m}\le2p_m^2e^{2p_m}\le4p_m^2.
\]
\end{proof}

\section{Four-direction algebra}
\begin{lemma}[Telescoping and rectangle identities]\label{lem:algebra}
For every ordering $i_1,\ldots,i_4$ of the directions,
\begin{equation}
 16F(x)-F(Dx)=\sum_{j=1}^4 2^{4-j}
 p_{i_j}(D_{i_{j-1}}\cdots D_{i_1}x).\label{eq:telescope}
\end{equation}
For distinct $i,j$,
\begin{equation}
 p_i(D_jx)=2p_i(x)-2p_j(x)+p_j(D_ix)
 \le2p_i(x)+p_j(D_ix).\label{eq:rectangle}
\end{equation}
Hence $D(r)\ge0$, and
\begin{equation}
 p_i(x^{[i]})\le D(r),\qquad
 x^{[i]}_i=r,\quad x^{[i]}_j=2r\ (j\ne i).\label{eq:opposite-face}
\end{equation}
\end{lemma}
\begin{proof}
Expand both equalities; intermediate values of $F$ cancel. The inequality
uses $p_j\ge0$. For \eqref{eq:opposite-face}, telescope from
$\underline r$ with direction $i$ last. The last term is
$p_i(x^{[i]})$ with coefficient one, and the other terms are nonnegative.
\end{proof}

Fix the order $1,2,3,4$ and write $\delta=D(r)$. Define
\begin{align*}
 y^{(i)}_j&=\begin{cases}4r&j<i,\\2r&j\ge i,\end{cases}
 &u_i&=p_i(y^{(i)}),\\
 z^{(i)}_j&=\begin{cases}4r&j<i,\\r&j=i,\\2r&j>i,\end{cases}
 &a_i&=p_i(z^{(i)}).
\end{align*}
Thus $u_i\le4a_i^2$ whenever $a_i\le1/4$, and
\begin{equation}
 D(2r)=8u_1+4u_2+2u_3+u_4.\label{eq:new-defect}
\end{equation}

\begin{lemma}[Preparing a coordinate before squaring]\label{lem:prepare}
For each $i=1,2,3,4$,
\begin{equation}
 a_i\le 2^{i-1}\delta+\sum_{j<i}2^{i-1-j}u_j.\label{eq:prepare}
\end{equation}
\end{lemma}
\begin{proof}
Start at $x^{[i]}$, where $p_i\le\delta$. Double the coordinates
$j=1,\ldots,i-1$ in order. At the step for $j$, the box $x$ has earlier
coordinates $4r$, coordinate $i$ equal to $r$, and the others equal to
$2r$. Thus $D_i x=y^{(j)}$, and the last term in
\eqref{eq:rectangle} is exactly $u_j$. Each step replaces a previous
bound $A$ by $2A+u_j$. Iterating proves \eqref{eq:prepare}.
No unproved monotonicity in transverse size is being used.
\end{proof}

\begin{proposition}[Quadratic improvement]\label{prop:recursion}
If $D(r)\le1/2048$, then
\begin{equation}
 D(2r)\le700D(r)^2\le1024D(r)^2.\label{eq:recursion}
\end{equation}
\end{proposition}
\begin{proof}
Use Lemmas~\ref{lem:prepare} and~\ref{lem:squaring} successively:
\[
\begin{array}{c|l|c|c}
 i&\text{unsquared bound}&a_i\text{ bounded by}&u_i\text{ bounded by}\\\hline
 1&\delta&\delta&4\delta^2\\
 2&2\delta+4\delta^2&3\delta&36\delta^2\\
 3&4\delta+44\delta^2&5\delta&100\delta^2\\
 4&8\delta+188\delta^2&9\delta&324\delta^2
\end{array}
\]
Indeed $4\delta,44\delta,188\delta\le1$ and
$9\delta\le9/2048<1/4$, so all squaring steps are legitimate.
Equation~\eqref{eq:new-defect} now gives
\[
 D(2r)\le(8\cdot4+4\cdot36+2\cdot100+324)\delta^2=700\delta^2.
\]
\end{proof}

\section{Iteration and the finite-cross-section gap}
\begin{proof}[Proof of Theorem~\ref{thm:main}]
Set $C=1024$, $\delta_k=D(2^kn)$, and $q=C\delta_0\le1/2$.
Equation~\eqref{eq:recursion} preserves the threshold. Induction gives
\begin{equation}
 \delta_k\le C^{-1}q^{2^k}\le C^{-1}2^{-2^k}.\label{eq:iterate}
\end{equation}
At the cube $\underline{w_k}$, telescope with the chosen time direction
first, so $8p_i(\underline{w_k})\le\delta_k$.
Lemma~\ref{lem:squaring} implies
\[
 (\lambda_2(w_k)/\lambda_1(w_k))^{w_k}
 \le2p_i(\underline{w_k})\le\delta_k/4
 \le(4C)^{-1}2^{-2^k}.
\]
Taking the $w_k=2^kn$ root proves the claim.
\end{proof}
The proof includes $D(n)=0$: the excited eigenvalue ratio is then zero.
This does not assert the existence of nontrivial finite-energy excitations.

\section{Local limits and the full observable Hilbert space}\label{sec:local}
\begin{assumption}[Local-observable interface]\label{ass:local}
The partition functions come from Borel probability laws on products of
compact metrizable local configuration spaces, invariant under translations
and the relevant time reflections. Normalize $\widehat K_w=K_w/\lambda_1(w)$ and its
leading vector $\Omega_w$. Every bounded local observable $A$ in a time
slab of thickness $r$ has an insertion $K_{w,A}$ satisfying
\[
 A_w:=K_{w,A}/\lambda_1(w)^r,\qquad\norm{A_w}\le\norm{A}_\infty,
\]
For torus periods $w>r$ containing the support, the torus and infinite-time
cylinder expectations are
\[
 \omega_{w,\mathrm{tor}}(A)=
 \frac{\Tr(A_w\widehat K_w^{w-r})}{\Tr\widehat K_w^w},\qquad
 \omega_{w,\mathrm{cyl}}(A)=\ip{\Omega_w}{A_w\Omega_w}.
\]
Separated insertions compose through powers of $\widehat K_w$.
The limiting states retain these invariances. On their admitted positive-time
local algebra, let $\Theta$ denote time reflection followed by complex
conjugation and set $\ip{[A]}{[B]}=\omega((\Theta A)B)$.
This form is positive semidefinite. The Hilbert space is its quotient
completion, with $\Omega=[1]$. The \emph{canonical one-step translation}
$T[A]=[\tau_1A]$ is well-defined on this quotient and extends to a positive
self-adjoint contraction. In particular $T\Omega=\Omega$ and
$\ip{[A]}{T^s[B]}=\omega((\Theta A)\tau_sB)$.
For a gauge theory the admitted algebra is gauge-invariant, the marginal
laws are gauge invariant, and the compact gauge group acts continuously;
the completion is the physical local-observable vacuum representation.
\end{assumption}
For a positive integral kernel, $|K_{w,A}|\le\norm{A}_\infty K_w^r$
implies the insertion norm: bound a matrix element using $|f|,|g|$ and
$\norm{K_w^r}=\lambda_1(w)^r$. The insertion need not itself be positive.

\begin{proposition}[Cylinder clustering and torus comparison]\label{prop:limits}
Under Assumptions~\ref{ass:trace} and~\ref{ass:local} and the small-test
hypothesis, put $\rho=2^{-1/n}$. Bounded observables $A,B$ separated by
$s$ transfer steps satisfy
\begin{equation}
 |\Cov_{\omega_{w_k,\mathrm{cyl}}}(A,B)|
 \le\norm{A}_\infty\norm{B}_\infty\rho^s.\label{eq:clustering}
\end{equation}
Every local subsequential limit of the cubic tori of side $w_k$ obeys the
same inequality. Its full physical transfer operator satisfies
$\norm{T|_{\Omega^\perp}}\le\rho$.
\end{proposition}
\begin{proof}
Insert $\widehat K_{w_k}^{s}-P_{\Omega_{w_k}}$ between $A_{w_k}$ and
$B_{w_k}$ and apply Theorem~\ref{thm:main} to obtain
\eqref{eq:clustering}.
First take continuous local observables. For a fixed local $A$, take
$k\ge1$ large enough that $r\le w_k/2$.
Set $s_m(w)=\Tr\widehat K_w^m-1$. Separating the leading eigenvector in
numerator and denominator shows
\begin{equation}
 |\omega_{w_k,\mathrm{tor}}(A)-\omega_{w_k,\mathrm{cyl}}(A)|
 \le2\norm{A}_\infty s_{w_k/2}(w_k).\label{eq:torus-cylinder}
\end{equation}
The opposite-face telescoping identity gives
\[
 p_1(w_k/2,w_k,w_k,w_k)\le D(w_k/2)=\delta_{k-1}.
\]
Thus $s_{w_k/2}(w_k)\le e^{\delta_{k-1}}-1\to0$.
Compactness supplies local weak subsequential limits; the comparison
identifies torus and cylinder limits along the same subsequence. Apply it
to continuous $A$, $B$, and $AB$, with fixed supports, to pass
\eqref{eq:clustering} to the limit. For bounded measurable local
observables, work in the limiting finite-support marginals: approximate
in $L^2$ by continuous functions with the same supremum bounds. Gauge
averaging preserves continuity, gauge invariance and these bounds.
Cauchy--Schwarz then passes the covariance bound to the measurable
observables. This is not a claim that weak convergence itself implies
convergence of every bounded measurable expectation.

A centered vector $v_A$ from a local positive-time observable has
$\ip{v_A}{T^sv_A}$ equal to a covariance of reflected slabs. Their finite
thickness changes only a prefactor, so
\[
 0\le\ip{v_A}{T^sv_A}\le C_A\rho^s
\]
for all sufficiently large integer $s$. Finite linear combinations retain
the same $\rho$. A positive spectral measure with nonzero mass in $[b,1]$
for some $b>\rho$ would give a lower bound proportional to $b^s$,
a contradiction. Bounded spectral projections and density extend the
exclusion to all of $\Omega^\perp$.
\end{proof}
This concerns each selected local-limit vacuum representation, not
uniqueness among all Gibbs states or external-charge superselection
sectors. It permits the vacuum complement to be zero.

\section{Wilson normalization and a strong-coupling check}\label{sec:wilson}
For pure four-dimensional $SU(2)$ Wilson theory, use normalized Haar per
positively oriented link and one plaquette per coordinate pair:
\begin{equation}
 Z_\beta(x)=\int\exp\left(\frac\beta2\sum_p\Tr U_p\right)
                 \prod_e dU_e,\qquad\beta\ge0.\label{eq:Wilson}
\end{equation}
Here $\Tr U_p$ is real and $|\Tr U_p/2|\le1$.
The convention $\beta\sum_p(1-\Tr U_p/2)$ differs by an extensive factor
that cancels in $D$. The structural transfer interface is recorded in
\texttt{audits/su2\_transfer\_interface.md}; it must establish the exact
normalized trace, gauge projection, operator positivity, and local insertions.

Once that interface is supplied, there is a crude direct strong-coupling
check. With all periods at least four, every plaquette has four distinct
link variables and Haar mean trace zero. There are
$N_p(x)=6\prod_i x_i$ plaquettes. Jensen's inequality and a pointwise bound give
\[
 0\le\log Z_\beta(x)\le6\beta\prod_i x_i,
 \qquad D_\beta(4)\le24576\beta.\label{eq:strong}
\]
For $0\le\beta\le2^{-26}$ the last bound is $3/8192<1/2048$.
The criterion therefore recovers a contraction at most $2^{-1/4}$ on
dyadic spatial cubes in this deliberately tiny strong-coupling interval,
provided the structural interface holds. This is a consistency check,
not a claimed advance in weak-coupling mass generation.

At $\beta=0$, $Z=1$, $D=0$, and the physical transfer operator is rank one.
The finite-slice gauge-invariant $L^2$ space is still infinite dimensional,
but all its orthogonal-to-vacuum directions are killed by transfer. The
reflected observable Hilbert space can therefore contain only the vacuum.
Nontriviality is a genuinely separate requirement.

\section{The exact missing continuum input}
An intended application would choose an asymptotically free bare-coupling
trajectory $\beta(a)$ and even $n(a)$ satisfying
\begin{equation}
 an(a)\to R_0\in(0,\infty),\qquad
 D_{\beta(a)}(n(a))\le1/2048
 \quad\text{for all sufficiently small }a.\label{eq:missing}
\end{equation}
Then the cutoff transfer masses would satisfy
\[
 \frac{m_{\mathrm{lat}}(a)}a\ge\frac{\log2}{an(a)}
 \longrightarrow\frac{\log2}{R_0}>0.
\]
Equation~\eqref{eq:missing} is \emph{not proved here}. Even if it were,
one would still need nontrivial limiting local Yang--Mills observables,
the reconstruction axioms, the required short-distance behavior, and
compatible convergence transferring the bound to their physical Hilbert
space. Convergence of normalized semigroup matrix elements on a dense
observable domain is one sufficient interface. A weak subsequential
measure limit or a cutoff gap alone does not suffice; compare \cite{clay}.

\begin{proposition}[What a cutoff-comparison lemma would suffice to prove]
Fix an integer blocking factor $B>1$ and let
$F_N(R)=\log Z_{\beta_N}(RB^N,RB^N,RB^N,RB^N)$ on admitted boxes.
Suppose a common coupling trajectory $\beta_N$ and volume-independent
bulk constants $c_N$ are chosen. Write
$\widetilde F_N(R)=F_N(R)-c_NR^4$ and
$\mathfrak D_N(R)=16F_N(R)-F_N(2R)$.
If, for one reference depth $K$, every $N\ge K$, and $R=M,2M$,
\[
 |\widetilde F_N(R)-\widetilde F_K(R)|\le C R^4\epsilon_K,
\]
then
\[
 |\mathfrak D_N(M)-\mathfrak D_K(M)|\le32CM^4\epsilon_K.
\]
In particular, $\mathfrak D_K(M)\le\varepsilon_*/2$ and
$32CM^4\epsilon_K\le\varepsilon_*/2$ imply a common successful physical
box of side $M$ at every finer cutoff. Under the structural transfer
hypotheses, the cutoff physical transfer masses are at least $\log2/M$.
\end{proposition}
\begin{proof}
The bulk terms cancel because $16M^4=(2M)^4$.
The triangle inequality bounds the two errors by
$16CM^4\epsilon_K+C(2M)^4\epsilon_K=32CM^4\epsilon_K$.
Apply Theorem~\ref{thm:main} with microscopic side $n_N=MB^N$ and spacing
$a_N=B^{-N}$.
\end{proof}
Both inputs to this proposition remain unproved for the required
Yang--Mills trajectory. The coupling and bulk subtraction must be common
to $M$ and $2M$; choosing a separate normalization for each box would
invalidate the conclusion. The elementary strong-coupling check in
Section~\ref{sec:wilson} does not supply a matched weak-coupling reference
at a depth where the comparison error is small.

The finite-box mechanism in \cite{repo-o3} balances rough mixing against
an error proportional to physical box volume. If a four-dimensional
analogue supplied
\[
 \ell_{\mathrm{mix}}(K)\lesssim B^{\gamma K},\qquad
 \mathrm{error}(K,R)\lesssim R^4 B^{-\alpha K},\quad B>1,
\]
then $R=B^{\delta K}$ would be both large enough to mix and small enough
to compare whenever $\gamma<\delta<\alpha/4$.
Thus $4\gamma<\alpha$ is sufficient for \emph{this particular strategy},
not necessary for Yang--Mills mass generation. Neither proposed
Yang--Mills estimate is established here. Mixing must control the
finite-volume defect or a sufficiently strong full-observable substitute;
one selected two-point correlation is not automatically enough.

\section{Verification and remaining scope}
\begin{itemize}
\item The Python check verifies exact integer coefficient identities for
all 24 telescoping orders and the rectangle identities, exact rational
scalar bounds, and high-precision regressions on positive spectra and
actual four-dimensional product-chain models. Examples do not replace proof.
\item The formal directory records the exact scope of any Lean certificate.
Rational constant certificates are not formalizations of the operator
theorem or of Yang--Mills theory.
\item Separate audits cover the source proof, the $SU(2)$ interface,
adversarial examples, cutoff obligations, and prior literature.
\item No published manuscript is modified. No claim is made that the full
Family 215 continuum construction has been independently verified, that
this criterion is new, or that the Millennium problem has been resolved.
\end{itemize}

\begin{thebibliography}{9}
\bibitem{repo-o4}
OpenAI, \emph{Sharp mass bounds for the two-dimensional $O(4)$ model},
September 23, 2026. Supplied repository, section ``A finite-volume
criterion for the full gap,'' source \texttt{build/sections/finite-volume.tex}.
\bibitem{repo-o3}
OpenAI, \emph{The canonical massive continuum limit of the two-dimensional
$O(3)$ model}, October 4, 2026. Supplied repository, section ``Trace
estimates uniform in the cutoff,'' source \texttt{build/sections/trace.tex}.
\bibitem{clay}
A. Jaffe and E. Witten, \emph{Quantum Yang--Mills Theory}, in
\emph{The Millennium Prize Problems}, Clay Mathematics Institute and AMS,
2006, pp. 129--152. Official description,
\url{https://www.claymath.org/wp-content/uploads/2022/06/yangmills.pdf}.
\end{thebibliography}
\end{document}
